\section{Empirical Investigation}
% BUDGET: ~2.5 pages incl. Figure 2 and Tables 3-4. FACT SOURCE: notes/report_sources.md Sec. 4.
% RULE: only TEST numbers are results. Dev numbers may appear once, as tuning context (e.g. lambda choice).

\subsection{Experimental Design}
\begin{table}[t]
\centering
\small
\begin{tabular}{p{1.6cm}p{5.2cm}}
\toprule
Factor & Levels \\
\midrule
Severity & 0, 1 word, 10, 25, 50, 75\% of words \\
Length band & 20--29, 30--44, 45--60 words \\
Methods & BERT + rerank; LLM few-shot; no repair (secondary: LLM + article) \\
Data & 150 test sentences, 750 rows per method; dev 20 (tuning only) \\
Metrics & BERTScore, \frr{} (primary); exact match, repair gain, format failures, contamination (secondary) \\
\bottomrule
\end{tabular}
\caption{Experimental design.}
\label{tab:design}
\end{table}
% Levels are nested within the same sentence, so all comparisons are paired (per sentence: better / same / worse).
% Tuning: lambda grid on dev only (lambda 8; FRR 0.042 -> 0.126 from lambda 1 to 8, flat after). Test split was opened per
%   method only after its config was frozen (a code gate), so no setting was chosen on test.
The design crosses three factors inside the same 150 test sentences: five damage
levels (1 word, then 10, 25, 50, 75\% of the sentence, plus a clean level 0), three
length bands (20--29, 30--44, 45--60 words), and three repair settings (BERT +
reranking, LLM few-shot, no repair), with LLM + full article as a secondary check.
All five damage levels reuse the same sentence and the same anchor word, so every
comparison across levels is paired: for one sentence we can ask directly whether
repair got better, worse, or stayed the same as damage grew, rather than comparing
averages over different sentences. The 20 dev sentences (disjoint from test, held
out from all reported numbers) were used only to pick settings before the test
split was scored: the reranking weight $\lambda$ was chosen on dev alone
(Section~\ref{sec:models}), and each method's setting was frozen by a code gate
before its test rows were touched, so nothing reported here was chosen by looking
at test results. One setting was still open at the time of writing, as already
flagged in Section~\ref{sec:models}: the LLM's exact model, host and prompt
version.

\subsection{Metrics}
\label{sec:metrics}
% BERTScore F1 of the repaired vs gold span \citep{zhang2020bertscore}: roberta-large, layer 17, baseline rescaling
%   (so values can go below 0). Fact Recovery Rate: share of gold fact tokens in the span restored exactly in their slot,
%   punctuation ignored. Exact words; CER = Levenshtein / gold span length; repair gain = CER before - CER after.
%   Format failure = wrong-length answer, scored as no repair. Contamination probe \citep{sainz2023contamination}:
%   first half of each clean test sentence, near-verbatim continuation = similarity >= 0.9.
% DEVIL'S ADVOCATE (SQ3): a wrong name can still get a high BERTScore; that is the reason for FRR.
%   Repair gain per row is dominated by short spans (one wrong word in a 1-word span gives CER > 1). Decide: pooled CER?
% OPTIONAL: split FRR into visible vs dropped fact slots (68 of 692 fact slots have no letters).
We score the repaired span, not the whole sentence, since only the span was
touched. BERTScore F1 \citep{zhang2020bertscore} compares the repaired span to the
gold span using \texttt{roberta-large} (layer 17), rescaled against a baseline so a
score can go below 0; this is our primary meaning-level metric. \frr{} is our
primary fact-level metric: the share of gold fact tokens in the span that the
repair restores exactly in their slot, punctuation ignored. As secondary metrics we
report exact-word match, character error rate (CER, a Levenshtein distance
normalised by gold span length) and repair gain (CER before minus CER after), the
LLM's format-failure rate, and a contamination probe \citep{sainz2023contamination}
that checks whether the LLM simply recites a memorised continuation of a clean test
sentence rather than genuinely repairing damaged ones.\footnote{Repair gain is
reported as a secondary number because it is dominated by short spans: one wrong
character in a one-word span already gives CER $>1$, so it does not aggregate
cleanly across levels; we do not pool it into a single figure.} A wrong-length LLM
answer is a format failure and is scored as no repair, never retried or hand-fixed.
BERTScore and \frr{} are both reported because they can disagree on purpose (SQ3):
a repair that substitutes a wrong but plausible name can still score well on
meaning while scoring zero on \frr{}, which is exactly the gap we want to measure.

\subsection{Results}
% Figure 2: level (x) vs span BERTScore and FRR (y), one line per method, no repair dashed. From notebook 04, TEST split.
% Table 3 below. Table 4 (optional): level x band, to separate damaged share from absolute gap length (SQ1).
\begin{table}[t]
\centering
\small\setlength{\tabcolsep}{4pt}
\begin{tabular}{lcccccc}
\toprule
 & \multicolumn{3}{c}{BERTScore (span)} & \multicolumn{3}{c}{\frr{}} \\
Level & BERT & LLM & None & BERT & LLM & None \\
\midrule
1w   & \pending & \pending & \pending & \pending & \pending & \pending \\
10\% & \pending & \pending & \pending & \pending & \pending & \pending \\
25\% & \pending & \pending & \pending & \pending & \pending & \pending \\
50\% & \pending & \pending & \pending & \pending & \pending & \pending \\
75\% & \pending & \pending & \pending & \pending & \pending & \pending \\
\bottomrule
\end{tabular}
\caption{Test results, 150 sentences per level. None = no repair. \todo{LLM format failures per level; paired counts.}}
\label{tab:results}
\end{table}
% Report in this order: (1) SQ1 curve shape and breakpoint; (2) band effect; (3) SQ2 paired BERT vs LLM per level;
%   (4) SQ3 BERTScore vs FRR gap; (5) LLM format-failure rate and where it occurs; (6) article vs few-shot (paired);
%   (7) contamination rate on test.
% WATCH (from dev, to check on test): on dev the methods crossed at 50-75% only because the LLM's format failures are
%   all there and are scored as no repair. If test shows the same, say the crossing is a format effect, not better repair.
\todo{Fill after the LLM test run and notebook 04 on test.}
